% --- Archon physics formalization source begin ---
% archon:physics
% archon:covers PhyXMiniProblems/problem_phyx_mini_0538.lean
% archon:source-report reports/phyx_mini/problem_phyx_mini_0538.source.json

\chapter{Physics problem phyx\_mini\_0538}
\label{ch:PhyXMiniProblems_problem_phyx_mini_0538}

\paragraph{Problem source.}
\#\# Physical scenario

Figure is a model of a benzene molecule. All atoms lie in a plane, and the carbon atoms (\textbackslash{}( m\_\{\textbackslash{}text\{C\}\} = 1.99 \textbackslash{}times 10\textasciicircum{}\{-26\} \textbackslash{}text\{ kg\} \textbackslash{})) form a regular hexagon, as do the hydrogen atoms (\textbackslash{}( m\_\{\textbackslash{}text\{H\}\} = 1.67 \textbackslash{}times 10\textasciicircum{}\{-27\} \textbackslash{}text\{ kg\} \textbackslash{})). The carbon atoms are 0.110 nm apart center to center, and the adjacent carbon and hydrogen atoms are 0.100 nm apart center to center.

\#\# Question

Calculate the moment of inertia of the molecule about an axis perpendicular to the plane of the paper through the center point \textbackslash{}( O \textbackslash{}).

\#\# Answer choices

- A: A: \textbackslash{}(2.31 \textbackslash{}times 10\textasciicircum{}\{-45\} \textbackslash{}text\{ kg\} \textbackslash{}cdot \textbackslash{}text\{m\}\textasciicircum{}2\textbackslash{})

- B: B: \textbackslash{}(1.65 \textbackslash{}times 10\textasciicircum{}\{-45\} \textbackslash{}text\{ kg\} \textbackslash{}cdot \textbackslash{}text\{m\}\textasciicircum{}2\textbackslash{})

- C: C: \textbackslash{}(1.05 \textbackslash{}times 10\textasciicircum{}\{-45\} \textbackslash{}text\{ kg\} \textbackslash{}cdot \textbackslash{}text\{m\}\textasciicircum{}2\textbackslash{})

- D: D: \textbackslash{}(1.89 \textbackslash{}times 10\textasciicircum{}\{-45\} \textbackslash{}text\{ kg\} \textbackslash{}cdot \textbackslash{}text\{m\}\textasciicircum{}2\textbackslash{})

\#\# Figure caption (auxiliary; use the image as primary evidence)

The image is a structural diagram of a benzene molecule. Here's a detailed description:

- **Molecular Structure**: The diagram represents a hexagonal ring structure, typical of benzene.

- **Atoms**: There are two types of atoms depicted:
  - **Carbons (C)**: Represented by orange circles. There are six carbon atoms forming the hexagonal ring.
  - **Hydrogens (H)**: Represented by blue circles, each attached to a carbon atom, extending outward.

- **Bonds**:
  - **Carbon-Carbon Bonds**: Lines connecting the carbon atoms in a hexagonal pattern.
  - **Carbon-Hydrogen Bonds**: Lines connecting each carbon to a hydrogen atom.

- **Measurements**:
  - The distance between carbon atoms is marked as 0.140 nm.
  - The distance from a carbon to a hydrogen is labeled as 0.110 nm.
  - The angle between the bonds at each carbon is noted to be 120°.

- **Color Legend**:
  - A color key indicates that blue circles represent hydrogen (H) and orange circles represent carbon (C).

These elements collectively describe the typical planar symmetry and bond angles found in benzene molecules.

\#\# Dataset metadata

Quantum Phenomena; Spatial Relation Reasoning, Physical Model Grounding Reasoning

\paragraph{Recorded answer/context.}
D: D: \textbackslash{}(1.89 \textbackslash{}times 10\textasciicircum{}\{-45\} \textbackslash{}text\{ kg\} \textbackslash{}cdot \textbackslash{}text\{m\}\textasciicircum{}2\textbackslash{})

\paragraph{Figure/image path.}
/root/proposal\_for\_physic/science-mango/phyx\_mini\_run/phyx\_data/test\_image/538.png

\paragraph{Formalization target.}
create a compiling Lean file with sorry bodies at `PhyXMiniProblems/problem\_phyx\_mini\_0538.lean`. The Lean declarations must preserve the physical quantities, dimensions or dimensional roles, figure labels, governing-law hypotheses, and final relation expressed by this problem.
Use Mathlib/Physlib names found through LeanExplore where available. If a physics API is missing, introduce faithful local abstractions rather than scalar placeholder aliases.

\begin{theorem}[Physics formalization target]
\label{thm:physics:phyx_mini_0538:target}
\lean{PhyXMiniProblems.ProblemPhyXMini0538.benzene_molecule_moment_of_inertia}
\uses{def:physics:phyx-mini-0538:phyxminiproblems-problemphyxmini0538-momentofinertiainkilogrammeterssquared, def:physics:phyx-mini-0538:phyxminiproblems-problemphyxmini0538-benzenemoleculesetup, def:physics:phyx-mini-0538:phyxminiproblems-problemphyxmini0538-matchesbenzenescenario, def:physics:phyx-mini-0538:phyxminiproblems-problemphyxmini0538-matchesproblemreadouts, def:physics:phyx-mini-0538:phyxminiproblems-problemphyxmini0538-matchesprimaryfigure, def:physics:phyx-mini-0538:phyxminiproblems-problemphyxmini0538-matchesplanarregularhexagongeometry, def:physics:phyx-mini-0538:phyxminiproblems-problemphyxmini0538-hasphysicalbenzeneparameters, def:physics:phyx-mini-0538:phyxminiproblems-problemphyxmini0538-satisfiespointmassmomentofinertialaw, def:physics:phyx-mini-0538:phyxminiproblems-problemphyxmini0538-answerchoice, def:physics:phyx-mini-0538:phyxminiproblems-problemphyxmini0538-displayedmomentofinertia, def:physics:phyx-mini-0538:phyxminiproblems-problemphyxmini0538-matchesanswerchoice}
The assigned autoformalize agent should translate this physics problem into Lean declarations in the covered file, with theorem and lemma proof bodies written as `by sorry`.
\end{theorem}
\begin{proof}
This is an autoformalization task, not a proof task. Produce statements that can later be proved by the physics prover without weakening the physical model.
\end{proof}

% --- Archon declaration topology begin ---
\section{Declaration topology}
\begin{definition}[moment Of Inertia Dimension]
  \label{def:physics:phyx-mini-0538:phyxminiproblems-problemphyxmini0538-momentofinertiadimension}
  \lean{PhyXMiniProblems.ProblemPhyXMini0538.momentOfInertiaDimension}
  The physical dimension of a moment of inertia, mass times length squared
\end{definition}
\begin{proof}
  This is the declared model definition of moment Of Inertia Dimension.
\end{proof}
\begin{definition}[Mass Quantity]
  \label{def:physics:phyx-mini-0538:phyxminiproblems-problemphyxmini0538-massquantity}
  \lean{PhyXMiniProblems.ProblemPhyXMini0538.MassQuantity}
  A nonnegative, unit-independent physical mass
\end{definition}
\begin{proof}
  This is the declared model definition of Mass Quantity.
\end{proof}
\begin{definition}[Length Quantity]
  \label{def:physics:phyx-mini-0538:phyxminiproblems-problemphyxmini0538-lengthquantity}
  \lean{PhyXMiniProblems.ProblemPhyXMini0538.LengthQuantity}
  A nonnegative, unit-independent physical length
\end{definition}
\begin{proof}
  This is the declared model definition of Length Quantity.
\end{proof}
\begin{definition}[Moment Of Inertia Quantity]
  \label{def:physics:phyx-mini-0538:phyxminiproblems-problemphyxmini0538-momentofinertiaquantity}
  \lean{PhyXMiniProblems.ProblemPhyXMini0538.MomentOfInertiaQuantity}
  \uses{def:physics:phyx-mini-0538:phyxminiproblems-problemphyxmini0538-momentofinertiadimension}
  A nonnegative physical moment of inertia about a specified axis
\end{definition}
\begin{proof}
  This is the declared model definition of Moment Of Inertia Quantity.
\end{proof}
\begin{definition}[mass Readout]
  \label{def:physics:phyx-mini-0538:phyxminiproblems-problemphyxmini0538-massreadout}
  \lean{PhyXMiniProblems.ProblemPhyXMini0538.massReadout}
  \uses{def:physics:phyx-mini-0538:phyxminiproblems-problemphyxmini0538-massquantity}
  Read a physical mass in a selected Physlib mass unit
\end{definition}
\begin{proof}
  This is the declared model definition of mass Readout.
\end{proof}
\begin{definition}[length Readout]
  \label{def:physics:phyx-mini-0538:phyxminiproblems-problemphyxmini0538-lengthreadout}
  \lean{PhyXMiniProblems.ProblemPhyXMini0538.lengthReadout}
  \uses{def:physics:phyx-mini-0538:phyxminiproblems-problemphyxmini0538-lengthquantity}
  Read a physical length in a selected Physlib length unit
\end{definition}
\begin{proof}
  This is the declared model definition of length Readout.
\end{proof}
\begin{definition}[moment Of Inertia Readout]
  \label{def:physics:phyx-mini-0538:phyxminiproblems-problemphyxmini0538-momentofinertiareadout}
  \lean{PhyXMiniProblems.ProblemPhyXMini0538.momentOfInertiaReadout}
  \uses{def:physics:phyx-mini-0538:phyxminiproblems-problemphyxmini0538-momentofinertiaquantity}
  Read a moment of inertia in coherent selected mass and length units
\end{definition}
\begin{proof}
  This is the declared model definition of moment Of Inertia Readout.
\end{proof}
\begin{definition}[mass In Kilograms]
  \label{def:physics:phyx-mini-0538:phyxminiproblems-problemphyxmini0538-massinkilograms}
  \lean{PhyXMiniProblems.ProblemPhyXMini0538.massInKilograms}
  \uses{def:physics:phyx-mini-0538:phyxminiproblems-problemphyxmini0538-massquantity, def:physics:phyx-mini-0538:phyxminiproblems-problemphyxmini0538-massreadout}
  Kilogram readout of an atom's mass
\end{definition}
\begin{proof}
  This is the declared model definition of mass In Kilograms.
\end{proof}
\begin{definition}[length In Meters]
  \label{def:physics:phyx-mini-0538:phyxminiproblems-problemphyxmini0538-lengthinmeters}
  \lean{PhyXMiniProblems.ProblemPhyXMini0538.lengthInMeters}
  \uses{def:physics:phyx-mini-0538:phyxminiproblems-problemphyxmini0538-lengthquantity, def:physics:phyx-mini-0538:phyxminiproblems-problemphyxmini0538-lengthreadout}
  Meter readout of a molecular length
\end{definition}
\begin{proof}
  This is the declared model definition of length In Meters.
\end{proof}
\begin{definition}[length In Nanometers]
  \label{def:physics:phyx-mini-0538:phyxminiproblems-problemphyxmini0538-lengthinnanometers}
  \lean{PhyXMiniProblems.ProblemPhyXMini0538.lengthInNanometers}
  \uses{def:physics:phyx-mini-0538:phyxminiproblems-problemphyxmini0538-lengthquantity, def:physics:phyx-mini-0538:phyxminiproblems-problemphyxmini0538-lengthreadout}
  Nanometer readout of a molecular length
\end{definition}
\begin{proof}
  This is the declared model definition of length In Nanometers.
\end{proof}
\begin{definition}[moment Of Inertia In Kilogram Meters Squared]
  \label{def:physics:phyx-mini-0538:phyxminiproblems-problemphyxmini0538-momentofinertiainkilogrammeterssquared}
  \lean{PhyXMiniProblems.ProblemPhyXMini0538.momentOfInertiaInKilogramMetersSquared}
  \uses{def:physics:phyx-mini-0538:phyxminiproblems-problemphyxmini0538-momentofinertiaquantity, def:physics:phyx-mini-0538:phyxminiproblems-problemphyxmini0538-momentofinertiareadout}
  SI readout of a moment of inertia, in kilogram-meter squared
\end{definition}
\begin{proof}
  This is the declared model definition of moment Of Inertia In Kilogram Meters Squared.
\end{proof}
\begin{definition}[Atom Kind]
  \label{def:physics:phyx-mini-0538:phyxminiproblems-problemphyxmini0538-atomkind}
  \lean{PhyXMiniProblems.ProblemPhyXMini0538.AtomKind}
  The two chemical species shown in the benzene diagram
\end{definition}
\begin{proof}
  This is the declared model definition of Atom Kind.
\end{proof}
\begin{definition}[Atom Site]
  \label{def:physics:phyx-mini-0538:phyxminiproblems-problemphyxmini0538-atomsite}
  \lean{PhyXMiniProblems.ProblemPhyXMini0538.AtomSite}
  \uses{def:physics:phyx-mini-0538:phyxminiproblems-problemphyxmini0538-atomkind}
  One of the six angular positions occupied by each species
\end{definition}
\begin{proof}
  This is the declared model definition of Atom Site.
\end{proof}
\begin{definition}[carbon Site]
  \label{def:physics:phyx-mini-0538:phyxminiproblems-problemphyxmini0538-carbonsite}
  \lean{PhyXMiniProblems.ProblemPhyXMini0538.carbonSite}
  \uses{def:physics:phyx-mini-0538:phyxminiproblems-problemphyxmini0538-atomsite}
  The carbon atom at angular position i
\end{definition}
\begin{proof}
  This is the declared model definition of carbon Site.
\end{proof}
\begin{definition}[hydrogen Site]
  \label{def:physics:phyx-mini-0538:phyxminiproblems-problemphyxmini0538-hydrogensite}
  \lean{PhyXMiniProblems.ProblemPhyXMini0538.hydrogenSite}
  \uses{def:physics:phyx-mini-0538:phyxminiproblems-problemphyxmini0538-atomsite}
  The hydrogen atom attached to the carbon at angular position i
\end{definition}
\begin{proof}
  This is the declared model definition of hydrogen Site.
\end{proof}
\begin{definition}[Atom Body Model]
  \label{def:physics:phyx-mini-0538:phyxminiproblems-problemphyxmini0538-atombodymodel}
  \lean{PhyXMiniProblems.ProblemPhyXMini0538.AtomBodyModel}
  The point-particle idealization used for every atom
\end{definition}
\begin{proof}
  This is the declared model definition of Atom Body Model.
\end{proof}
\begin{definition}[Ring Shape]
  \label{def:physics:phyx-mini-0538:phyxminiproblems-problemphyxmini0538-ringshape}
  \lean{PhyXMiniProblems.ProblemPhyXMini0538.RingShape}
  The qualitative shape of each species' six-atom ring
\end{definition}
\begin{proof}
  This is the declared model definition of Ring Shape.
\end{proof}
\begin{definition}[Rotation Axis Geometry]
  \label{def:physics:phyx-mini-0538:phyxminiproblems-problemphyxmini0538-rotationaxisgeometry}
  \lean{PhyXMiniProblems.ProblemPhyXMini0538.RotationAxisGeometry}
  Geometry of the axis requested in the question
\end{definition}
\begin{proof}
  This is the declared model definition of Rotation Axis Geometry.
\end{proof}
\begin{definition}[Bond Kind]
  \label{def:physics:phyx-mini-0538:phyxminiproblems-problemphyxmini0538-bondkind}
  \lean{PhyXMiniProblems.ProblemPhyXMini0538.BondKind}
  The two kinds of labeled bonds in the primary image
\end{definition}
\begin{proof}
  This is the declared model definition of Bond Kind.
\end{proof}
\begin{definition}[Figure Color]
  \label{def:physics:phyx-mini-0538:phyxminiproblems-problemphyxmini0538-figurecolor}
  \lean{PhyXMiniProblems.ProblemPhyXMini0538.FigureColor}
  Colors used by the atom legend in the primary image
\end{definition}
\begin{proof}
  This is the declared model definition of Figure Color.
\end{proof}
\begin{definition}[Planar Coordinate Meters]
  \label{def:physics:phyx-mini-0538:phyxminiproblems-problemphyxmini0538-planarcoordinatemeters}
  \lean{PhyXMiniProblems.ProblemPhyXMini0538.PlanarCoordinateMeters}
  A scalar coordinate pair whose two components are explicitly in meters
\end{definition}
\begin{proof}
  This is the declared model definition of Planar Coordinate Meters.
\end{proof}
\begin{definition}[regular Hexagon Vertex Meters]
  \label{def:physics:phyx-mini-0538:phyxminiproblems-problemphyxmini0538-regularhexagonvertexmeters}
  \lean{PhyXMiniProblems.ProblemPhyXMini0538.regularHexagonVertexMeters}
  \uses{def:physics:phyx-mini-0538:phyxminiproblems-problemphyxmini0538-planarcoordinatemeters}
  Vertex i of a regular planar hexagon with a common angular phase
\end{definition}
\begin{proof}
  This is the declared model definition of regular Hexagon Vertex Meters.
\end{proof}
\begin{definition}[planar Distance Meters]
  \label{def:physics:phyx-mini-0538:phyxminiproblems-problemphyxmini0538-planardistancemeters}
  \lean{PhyXMiniProblems.ProblemPhyXMini0538.planarDistanceMeters}
  \uses{def:physics:phyx-mini-0538:phyxminiproblems-problemphyxmini0538-planarcoordinatemeters}
  Euclidean distance between two explicitly meter-valued planar coordinates
\end{definition}
\begin{proof}
  This is the declared model definition of planar Distance Meters.
\end{proof}
\begin{definition}[Benzene Figure]
  \label{def:physics:phyx-mini-0538:phyxminiproblems-problemphyxmini0538-benzenefigure}
  \lean{PhyXMiniProblems.ProblemPhyXMini0538.BenzeneFigure}
  \uses{def:physics:phyx-mini-0538:phyxminiproblems-problemphyxmini0538-atomkind, def:physics:phyx-mini-0538:phyxminiproblems-problemphyxmini0538-atomsite, def:physics:phyx-mini-0538:phyxminiproblems-problemphyxmini0538-bondkind, def:physics:phyx-mini-0538:phyxminiproblems-problemphyxmini0538-figurecolor}
  Presentation-level data visible in the supplied image. Quantitative physical lengths remain fields of the molecule setup; the figure stores only its printed scalar readouts and labels
\end{definition}
\begin{proof}
  This is the declared model definition of Benzene Figure.
\end{proof}
\begin{definition}[Benzene Molecule Setup]
  \label{def:physics:phyx-mini-0538:phyxminiproblems-problemphyxmini0538-benzenemoleculesetup}
  \lean{PhyXMiniProblems.ProblemPhyXMini0538.BenzeneMoleculeSetup}
  \uses{def:physics:phyx-mini-0538:phyxminiproblems-problemphyxmini0538-massquantity, def:physics:phyx-mini-0538:phyxminiproblems-problemphyxmini0538-lengthquantity, def:physics:phyx-mini-0538:phyxminiproblems-problemphyxmini0538-momentofinertiaquantity, def:physics:phyx-mini-0538:phyxminiproblems-problemphyxmini0538-atomkind, def:physics:phyx-mini-0538:phyxminiproblems-problemphyxmini0538-atomsite, def:physics:phyx-mini-0538:phyxminiproblems-problemphyxmini0538-atombodymodel, def:physics:phyx-mini-0538:phyxminiproblems-problemphyxmini0538-ringshape, def:physics:phyx-mini-0538:phyxminiproblems-problemphyxmini0538-rotationaxisgeometry, def:physics:phyx-mini-0538:phyxminiproblems-problemphyxmini0538-planarcoordinatemeters, def:physics:phyx-mini-0538:phyxminiproblems-problemphyxmini0538-benzenefigure}
  The physical molecule and its observables. The planar coordinates explicitly encode coplanarity. Neither the axis distances nor the moment of inertia are defined from the requested answer; their geometric and dynamical relations are imposed separately below
\end{definition}
\begin{proof}
  This is the declared model definition of Benzene Molecule Setup.
\end{proof}
\begin{definition}[Matches Benzene Scenario]
  \label{def:physics:phyx-mini-0538:phyxminiproblems-problemphyxmini0538-matchesbenzenescenario}
  \lean{PhyXMiniProblems.ProblemPhyXMini0538.MatchesBenzeneScenario}
  \uses{def:physics:phyx-mini-0538:phyxminiproblems-problemphyxmini0538-benzenemoleculesetup}
  Qualitative molecular and rotation-axis model stated in the problem
\end{definition}
\begin{proof}
  This is the declared model definition of Matches Benzene Scenario.
\end{proof}
\begin{definition}[Matches Problem Readouts]
  \label{def:physics:phyx-mini-0538:phyxminiproblems-problemphyxmini0538-matchesproblemreadouts}
  \lean{PhyXMiniProblems.ProblemPhyXMini0538.MatchesProblemReadouts}
  \uses{def:physics:phyx-mini-0538:phyxminiproblems-problemphyxmini0538-massinkilograms, def:physics:phyx-mini-0538:phyxminiproblems-problemphyxmini0538-lengthinnanometers, def:physics:phyx-mini-0538:phyxminiproblems-problemphyxmini0538-benzenemoleculesetup}
  Numerical data stated in the problem: carbon and hydrogen atom masses and the two center-to-center separations. No moment-of-inertia value occurs here
\end{definition}
\begin{proof}
  This is the declared model definition of Matches Problem Readouts.
\end{proof}
\begin{definition}[Matches Primary Figure]
  \label{def:physics:phyx-mini-0538:phyxminiproblems-problemphyxmini0538-matchesprimaryfigure}
  \lean{PhyXMiniProblems.ProblemPhyXMini0538.MatchesPrimaryFigure}
  \uses{def:physics:phyx-mini-0538:phyxminiproblems-problemphyxmini0538-lengthinnanometers, def:physics:phyx-mini-0538:phyxminiproblems-problemphyxmini0538-atomsite, def:physics:phyx-mini-0538:phyxminiproblems-problemphyxmini0538-bondkind, def:physics:phyx-mini-0538:phyxminiproblems-problemphyxmini0538-benzenemoleculesetup}
  Literal readouts from the primary image: six atoms of each species, six bonds of each kind, the blue H and orange C legend, center label O, distance labels 0.110 nm and 0.100 nm, and the 120 degree bond-angle label
\end{definition}
\begin{proof}
  This is the declared model definition of Matches Primary Figure.
\end{proof}
\begin{definition}[Matches Planar Regular Hexagon Geometry]
  \label{def:physics:phyx-mini-0538:phyxminiproblems-problemphyxmini0538-matchesplanarregularhexagongeometry}
  \lean{PhyXMiniProblems.ProblemPhyXMini0538.MatchesPlanarRegularHexagonGeometry}
  \uses{def:physics:phyx-mini-0538:phyxminiproblems-problemphyxmini0538-lengthreadout, def:physics:phyx-mini-0538:phyxminiproblems-problemphyxmini0538-lengthinmeters, def:physics:phyx-mini-0538:phyxminiproblems-problemphyxmini0538-carbonsite, def:physics:phyx-mini-0538:phyxminiproblems-problemphyxmini0538-hydrogensite, def:physics:phyx-mini-0538:phyxminiproblems-problemphyxmini0538-regularhexagonvertexmeters, def:physics:phyx-mini-0538:phyxminiproblems-problemphyxmini0538-planardistancemeters, def:physics:phyx-mini-0538:phyxminiproblems-problemphyxmini0538-benzenemoleculesetup}
  The two species occupy concentric regular hexagons with a common angular phase. For a regular hexagon the carbon ring's radius equals its side length; each attached hydrogen lies radially outward by one C-H separation. The last two fields state the corresponding perpendicular distances to the normal axis in every coherent length unit. These are geometric inputs, not inertia data
\end{definition}
\begin{proof}
  This is the declared model definition of Matches Planar Regular Hexagon Geometry.
\end{proof}
\begin{definition}[Has Physical Benzene Parameters]
  \label{def:physics:phyx-mini-0538:phyxminiproblems-problemphyxmini0538-hasphysicalbenzeneparameters}
  \lean{PhyXMiniProblems.ProblemPhyXMini0538.HasPhysicalBenzeneParameters}
  \uses{def:physics:phyx-mini-0538:phyxminiproblems-problemphyxmini0538-massinkilograms, def:physics:phyx-mini-0538:phyxminiproblems-problemphyxmini0538-lengthinmeters, def:physics:phyx-mini-0538:phyxminiproblems-problemphyxmini0538-momentofinertiainkilogrammeterssquared, def:physics:phyx-mini-0538:phyxminiproblems-problemphyxmini0538-atomsite, def:physics:phyx-mini-0538:phyxminiproblems-problemphyxmini0538-benzenemoleculesetup}
  Positivity and nondegeneracy conditions for the physical molecule
\end{definition}
\begin{proof}
  This is the declared model definition of Has Physical Benzene Parameters.
\end{proof}
\begin{definition}[Satisfies Point Mass Moment Of Inertia Law]
  \label{def:physics:phyx-mini-0538:phyxminiproblems-problemphyxmini0538-satisfiespointmassmomentofinertialaw}
  \lean{PhyXMiniProblems.ProblemPhyXMini0538.SatisfiesPointMassMomentOfInertiaLaw}
  \uses{def:physics:phyx-mini-0538:phyxminiproblems-problemphyxmini0538-massreadout, def:physics:phyx-mini-0538:phyxminiproblems-problemphyxmini0538-lengthreadout, def:physics:phyx-mini-0538:phyxminiproblems-problemphyxmini0538-momentofinertiareadout, def:physics:phyx-mini-0538:phyxminiproblems-problemphyxmini0538-atomsite, def:physics:phyx-mini-0538:phyxminiproblems-problemphyxmini0538-benzenemoleculesetup}
  The general point-mass law I = sum m\_i r\_i\textasciicircum{}2, stated in every coherent mass and length unit. It contains neither the benzene-specific simplified formula nor any displayed numerical answer
\end{definition}
\begin{proof}
  This is the declared model definition of Satisfies Point Mass Moment Of Inertia Law.
\end{proof}
\begin{definition}[Answer Choice]
  \label{def:physics:phyx-mini-0538:phyxminiproblems-problemphyxmini0538-answerchoice}
  \lean{PhyXMiniProblems.ProblemPhyXMini0538.AnswerChoice}
  Labels of the four answer choices in the source problem
\end{definition}
\begin{proof}
  This is the declared model definition of Answer Choice.
\end{proof}
\begin{definition}[displayed Moment Of Inertia]
  \label{def:physics:phyx-mini-0538:phyxminiproblems-problemphyxmini0538-displayedmomentofinertia}
  \lean{PhyXMiniProblems.ProblemPhyXMini0538.displayedMomentOfInertia}
  \uses{def:physics:phyx-mini-0538:phyxminiproblems-problemphyxmini0538-answerchoice}
  Displayed moments of inertia, read in kilogram-meter squared
\end{definition}
\begin{proof}
  This is the declared model definition of displayed Moment Of Inertia.
\end{proof}
\begin{definition}[displayed Resolution]
  \label{def:physics:phyx-mini-0538:phyxminiproblems-problemphyxmini0538-displayedresolution}
  \lean{PhyXMiniProblems.ProblemPhyXMini0538.displayedResolution}
  One unit in the final displayed digit of every answer choice
\end{definition}
\begin{proof}
  This is the declared model definition of displayed Resolution.
\end{proof}
\begin{definition}[Matches Answer Choice]
  \label{def:physics:phyx-mini-0538:phyxminiproblems-problemphyxmini0538-matchesanswerchoice}
  \lean{PhyXMiniProblems.ProblemPhyXMini0538.MatchesAnswerChoice}
  \uses{def:physics:phyx-mini-0538:phyxminiproblems-problemphyxmini0538-momentofinertiainkilogrammeterssquared, def:physics:phyx-mini-0538:phyxminiproblems-problemphyxmini0538-benzenemoleculesetup, def:physics:phyx-mini-0538:phyxminiproblems-problemphyxmini0538-answerchoice, def:physics:phyx-mini-0538:phyxminiproblems-problemphyxmini0538-displayedmomentofinertia, def:physics:phyx-mini-0538:phyxminiproblems-problemphyxmini0538-displayedresolution}
  Rounding agreement with a displayed answer choice
\end{definition}
\begin{proof}
  This is the declared model definition of Matches Answer Choice.
\end{proof}
\begin{lemma}[carbon axis radius in meters]
  \label{lem:physics:phyx-mini-0538:phyxminiproblems-problemphyxmini0538-carbon-axis-radius-in-meters}
  \lean{PhyXMiniProblems.ProblemPhyXMini0538.carbon_axis_radius_in_meters}
  \uses{def:physics:phyx-mini-0538:phyxminiproblems-problemphyxmini0538-lengthinmeters, def:physics:phyx-mini-0538:phyxminiproblems-problemphyxmini0538-carbonsite, def:physics:phyx-mini-0538:phyxminiproblems-problemphyxmini0538-benzenemoleculesetup, def:physics:phyx-mini-0538:phyxminiproblems-problemphyxmini0538-matchesplanarregularhexagongeometry}
  Every carbon lies one C-C side length from the normal axis through O
\end{lemma}
\begin{proof}
  The conclusion follows from the governing relations and figure readouts named in the statement of carbon axis radius in meters.
\end{proof}
\begin{lemma}[hydrogen axis radius in meters]
  \label{lem:physics:phyx-mini-0538:phyxminiproblems-problemphyxmini0538-hydrogen-axis-radius-in-meters}
  \lean{PhyXMiniProblems.ProblemPhyXMini0538.hydrogen_axis_radius_in_meters}
  \uses{def:physics:phyx-mini-0538:phyxminiproblems-problemphyxmini0538-lengthinmeters, def:physics:phyx-mini-0538:phyxminiproblems-problemphyxmini0538-hydrogensite, def:physics:phyx-mini-0538:phyxminiproblems-problemphyxmini0538-benzenemoleculesetup, def:physics:phyx-mini-0538:phyxminiproblems-problemphyxmini0538-matchesplanarregularhexagongeometry}
  Every hydrogen radius is the carbon radius plus one radial C-H bond
\end{lemma}
\begin{proof}
  The conclusion follows from the governing relations and figure readouts named in the statement of hydrogen axis radius in meters.
\end{proof}
\begin{lemma}[benzene point mass inertia formula]
  \label{lem:physics:phyx-mini-0538:phyxminiproblems-problemphyxmini0538-benzene-point-mass-inertia-formula}
  \lean{PhyXMiniProblems.ProblemPhyXMini0538.benzene_point_mass_inertia_formula}
  \uses{def:physics:phyx-mini-0538:phyxminiproblems-problemphyxmini0538-massinkilograms, def:physics:phyx-mini-0538:phyxminiproblems-problemphyxmini0538-lengthinmeters, def:physics:phyx-mini-0538:phyxminiproblems-problemphyxmini0538-momentofinertiainkilogrammeterssquared, def:physics:phyx-mini-0538:phyxminiproblems-problemphyxmini0538-benzenemoleculesetup, def:physics:phyx-mini-0538:phyxminiproblems-problemphyxmini0538-matchesplanarregularhexagongeometry, def:physics:phyx-mini-0538:phyxminiproblems-problemphyxmini0538-satisfiespointmassmomentofinertialaw}
  Grouping the twelve terms of the general point-mass sum by species gives six equal carbon contributions and six equal hydrogen contributions. This formula is derived from geometry and the general law; it is not a premise
\end{lemma}
\begin{proof}
  The conclusion follows from the governing relations and figure readouts named in the statement of benzene point mass inertia formula.
\end{proof}
% --- Archon declaration topology end ---
% --- Archon physics formalization source end ---
